\section{Methodology}
\label{sec:method}

% ==================== LR START ====================

\subsection{Dataset}
\label{subsec:dataset}

In this study, we used the \texttt{RQ1} subset of the
\textit{ImproBR Replication Package}, which is publicly available on
Figshare~\cite{improbr2026dataset}. It contains 139 real-world bug
reports collected from Mojira, the official issue tracker for
Minecraft. We selected \texttt{RQ1} because it provides the report
content and manual labels required to evaluate the executability of
\textit{Steps to Reproduce} (S2R).

Each issue has two corresponding versions. The \textit{Raw} version
contains the original report submitted to Mojira before being processed
by ImproBR. The \textit{Improved} version is the corresponding output
of the ImproBR pipeline. In the Improved report, components such as S2R
may be added, clarified, or reorganized. The two versions were paired
using \texttt{issue\_key}, enabling a paired comparison of the same
issue under the Raw and Improved conditions.

Ground-truth labels were obtained from the manual evaluations included
in the replication package~\cite{improbr2026dataset}. Two independent
human annotators evaluated both report versions. Because our study
focuses on automatically assessing S2R executability, we used only the
S2R label. This label has two possible values:
\texttt{Executable} and \texttt{Non-Executable}. A report is labeled
\texttt{Executable} when its reproduction steps are sufficiently clear
for an evaluator to follow. A report is labeled
\texttt{Non-Executable} when missing, incorrect, or ambiguous
information prevents the evaluator from following the steps without
making additional assumptions.

The Raw condition contains 40 \emph{Executable} reports (28.8\%) and
99 \emph{Non-Executable} reports (71.2\%). The Improved condition
contains 94 \emph{Executable} reports (67.6\%) and 45
\emph{Non-Executable} reports (32.4\%). These ground-truth labels serve
as the reference for evaluating model predictions. Each prediction is
matched to its corresponding label using \texttt{issue\_key} and the
report condition (Raw or Improved).

The 139 issues were divided into three subsets:
\textit{Pilot}, \textit{Development}, and \textit{Holdout}. First,
26 issues were assigned to Pilot for checking the execution pipeline,
output schema, and prompt design. After excluding the Pilot issues, the
remaining 113 issues were divided into 75 Development issues and 38
Holdout issues. We used \texttt{seed}=210 to make this split
reproducible.

The Development set was used to refine the Full configuration and
conduct error analysis; its results were treated as calibration
evidence rather than a final independent performance estimate. The
Holdout set was reserved for the final evaluation. Splitting was
performed at the \texttt{issue\_key} level, so the Raw and Improved
versions of each issue were always assigned to the same subset. This
prevents pair-level leakage between the Development and Holdout sets.


\subsection{Experimental Setup and Pipeline}
\label{subsec:pipeline}

The experimental pipeline was implemented in Python and accessed the
OpenAI Chat Completions API through the \texttt{openai} SDK. We used
\texttt{gpt-4o-mini-2024-07-18} with
\texttt{temperature}=0 and \texttt{seed}=210. JSON mode was enabled
using \path|response_format={"type":"json_object"}|. The
\texttt{max\_tokens} parameter was not explicitly specified. Fixing
the model version and generation parameters reduces unnecessary output
variation and supports reproducibility, although it does not guarantee
identical responses across repeated executions.

We used separate prompt templates for the Raw and Improved conditions.
Each template consists of System and User sections. The System section
defines the model's role, evaluation rubric, allowed label values, and
JSON output requirement. The User section supplies the
\texttt{issue\_key}, report content, and required JSON schema. The Raw
prompt instructs the model to evaluate only the Raw report. The
Improved prompt specifies the Improved report as the evaluation target
and includes its paired Raw report only as context for preserving the
original issue intent. Owing to the page limit, the complete prompt
texts are provided in the replication package.

During the Pilot phase, we used the Raw V10 and Improved V18 prompts,
both based on the rubric defined in
\path{evaluation_metrics.yaml}. For the Full phase, Raw V11 and
Improved V20 were refined on the Development set. Improved V20 combines
the V19 prompt with Rules V20. The resulting Full configurations were
then evaluated on the same 26 Pilot issues during Pilot Validation.

Post-processing rules were defined according to the evaluation rubric.
Pilot-specific rules were refined during the Pilot phase, whereas rules
for the Full configuration were calibrated only on the Development
set. During inference, these rules inspect the target report and, where
applicable in the Improved condition, its paired Raw context.
Ground Truth and prior predictions are not available to the rules.
Holdout labels and mismatches were not used for prompt refinement or
rule calibration.

After Pilot Validation, we froze the model version, prompts,
experimental script, post-processing rules, and metric logic. We
recorded SHA-256 hashes of the frozen files to verify their integrity
before the Holdout evaluation. The frozen configuration was then
evaluated once on the 38 Holdout issues for each condition. Holdout
results were not used for further tuning.

For each issue, the pipeline reads the common fields
\emph{Issue Key}, \emph{Summary}, \emph{Type},
\emph{Affects Version/s}, \emph{Labels}, \emph{Category}, and
\emph{Description}. Improved reports additionally include
\emph{Steps to Reproduce}, \emph{Observed Behavior},
\emph{Expected Behavior}, and \emph{Environment}. Ground Truth,
annotator labels, prior predictions, and potentially biasing
administrative fields---including \emph{Confirmation Status},
\emph{Resolution}, \emph{Fix Version/s}, \emph{Votes}, and
\emph{Watch Count}---are excluded from the prompt.

For the Raw condition, the Raw report is inserted into the
corresponding Raw prompt. For the Improved condition, the Improved
report is the evaluation target, while its paired Raw report is
included only as original-intent context. Each issue in each condition
is processed through one application-level API request. Ground Truth
is not included in the request and is introduced only during metric
computation.

The model is required to return a structured JSON object containing
the \texttt{issue\_key}; the S2R label, reproducibility, validity,
failure type, and rationale; auxiliary assessments of Expected and
Observed Behavior; and an overall rationale. The
\path{run_experiment.py} script reads the CSV input, identifies the
\texttt{issue\_key}, constructs the prompt, calls the API, parses the
response, applies the post-processing rules, and stores the result.

After receiving a response, the pipeline checks whether it is valid
JSON, verifies that all required fields are present, normalizes allowed
label values, and checks logical consistency among the fields. If an
API error, empty response, malformed JSON, or missing required field
occurs, the script records the case with status \texttt{error} and
continues with the next issue. The experiment script does not implement
an additional retry loop.

Structured predictions are stored in
\path{*_llm_output_*.csv}, while API status, model and prompt versions,
seed, \texttt{system\_fingerprint}, token usage, and cost are stored in
\path{*_api_log_*.csv}. Predictions are then matched with Ground Truth
using both \texttt{issue\_key} and the corresponding condition
(Raw or Improved). The merged data are passed to
\path{compute_metric.py}, which computes Accuracy, Cohen's Kappa, the
confusion matrix, and the list of mismatched cases. The source code,
prompts, data splits, and raw experimental outputs required for
replication are available at
\url{https://github.com/LamSongToan/SWT301_SE1925_G6}.

% ===================== LR END =====================
% ===================== MS START ===================

% MS writes Metrics and Statistical Analysis Plan below.
